\chapter*{Quy ước cập nhật báo cáo}
\addcontentsline{toc}{chapter}{Quy ước cập nhật báo cáo}
Báo cáo này được cả nhóm viết dần trong quá trình phát triển, để thành viên mới đọc hiểu được project và để báo cáo
cuối kỳ không phải soạn lại từ đầu. Quy ước:
\begin{itemize}
  \item Phần chưa làm chỉ có đề mục, kèm dòng \textit{(Dự kiến, chưa thực hiện.)} -- macro \code{\textbackslash chualam}.
  \item Phần đã viết nhưng còn sửa theo các tuần sau có ghi chú nhỏ bên dưới -- macro
        \code{\textbackslash dukiensua\{nội dung, tuần\}}.
  \item Thông tin còn thiếu đánh dấu màu đỏ -- macro \code{\textbackslash canbosung}. Trên trang bìa, mỗi thành viên tự
        điền họ tên và MSSV của mình; trưởng nhóm điền học phần và giảng viên.
  \item Mỗi việc trong \code{docs/checkpoints.md} khi chuyển sang DONE thì owner cập nhật mục tương ứng ở đây trong cùng
        PR, kèm số liệu và bằng chứng tái hiện được. Không ghi kết quả chưa đo.
  \item Nguồn sự thật cho giao diện giữa các module vẫn là \code{docs/contracts.md}; báo cáo giải thích và dẫn lại, không
        định nghĩa contract mới.
\end{itemize}

\begin{table}[H]
\centering\small
\renewcommand{\arraystretch}{1.18}
\caption{Lịch sử cập nhật báo cáo}\label{tab:lichsu}
\begin{tabularx}{\linewidth}{@{}llX@{}}
\toprule
\textbf{Ngày} & \textbf{Người} & \textbf{Nội dung} \\ \midrule
15/10/2026 & Nhân & Khởi tạo khung báo cáo; viết phần CP0 (kiến trúc, luồng dữ liệu, serving, kiểm thử) \\
15/10/2026 & Nhân & Thêm use case (biểu đồ, đặc tả), biểu đồ hoạt động; viết lại Chương 2 kèm hình minh họa; chuyển
tên file, lệnh sang Phụ lục A \\
23/11/2026 & Nhân & Lược đồ năm bảng mới, API ba nghiệp vụ và định dạng lỗi chung, đặc tả ba use case; kết quả trên dữ
liệu mẫu, retry, dung lượng và hạn lưu 7 ngày (W2-NHAN-01, W3-NHAN-01) \\
23/11/2026 & Nhân & Tầng kiểm thử hồi quy contract, kết quả hồi quy sau khi gộp, phụ lục mô tả API và hướng dẫn vận
hành (W4-NHAN-01) \\
\bottomrule
\end{tabularx}
\end{table}
